\chapter{Conclusion}
\label{ch:conclusion}

A genomic resistance model can look excellent on paper and still be measuring the
wrong thing. That doubt is where this research started. The question we set out to
answer was how much of the score comes from real biological patterns, and how much
comes from similarities between related strains. We had to have a powerful model in
order to test this fairly. TriLens-MoE-Stack was able to obtain an average PR-AUC of 0.849 on the random
split, the best score of all the models tested. Against every method from
outside this work the lead is clear: 0.045 over the PCA gradient-boosting model this
project began with, 0.154 over TabNet, and 0.557 over the null model, and on three of the seven
antibiotics the score clears the upper bound of that baseline's confidence interval.
Against its own strongest component the margin is 0.009, which is what an ablation is
expected to show once the best part is already in place. Then whole sequence types were removed from
training and no test strain had a close neighbour in the training set. The average
PR-AUC decreased to 0.634. A portion of this decrease was due to a decreased fraction
of resistant strains in the holdout set which has a direct impact on PR-AUC. This
change was corrected, and the score was 0.723. The actual difference was not 0.215
but 0.126. Also, the decline in performance was quite distinct across antibiotics.
Three antibiotics lost only 0.017 to 0.031, while four lost 0.176 to 0.244. This
stark division indicates that the resistance associated with the stable biological
processes can be transferred to new lineages, whereas the patterns primarily
depending on the ancestry are not very transferable. When the classical part of the
pipeline was rebuilt on a second dataset, it kept about 90\% of its earlier
performance. This indicates that our approach is not just fitted to the specific
characteristics of a single dataset.

However, there are also significant limitations to these results. None of the
antibiotics passed both CLSI error limits, so the error rates stay above the level
those limits set. That weakness is concentrated in a few drugs rather than spread
across the panel. For ciprofloxacin and ampicillin the model reaches 0.9773 and
0.9440 on the random split, and still holds 0.9118 and 0.9043 when every close
relative is removed from training. The lineage holdout led to a higher share of resistant strains being incorrectly
classified as susceptible (0.249 to 0.559), although the accuracy-like measure did
not change significantly. This distinction is important since the accuracy-like
measure is commonly reported in the majority of studies and it may conceal severe
clinical errors. The results of the interpretation were also not entirely consistent.
In the attention weights, only one known resistance gene was identified whereas eight
were identified using post-hoc attribution to classical learners. Thus the gene-level
results should be viewed as potential leads rather than established biological
truths. There is also a definite limit to the gene presence and absence lens. It is
able to identify the presence of a gene, but not the level of gene expression. This
may account for amoxicillin-clavulanate being the most difficult antibiotic during
the study. One more limitation arises from the stratified split that is based on the
phylogroup labels. These labels were extracted from an unreferenced lookup table for
57 sequence types and we were not able to verify the original source of that table.

The key takeaways from this work are related to testing and refining genomic
resistance models. An adequate lineage holdout should eliminate entire clonal groups,
but it should also guard against the domination of the test set by a single large
clone. This produces another issue as the prevalence of resistance can be reduced
with the removal of large resistant groups, which can then lead to a reduction in
PR-AUC even if the true ability of the model has not reduced as much. In this study,
prevalence correction moved the generalisation gap from 0.215 to 0.126. Thus,
approximately two fifths of the original drop was caused by the evaluation measure
instead of an actual loss in model performance. A critical issue with gated attention
pooling for binary gene data was also uncovered in this study. The original attention
score could be the same for all strains, as each present gene used only its fixed
embedding. When the score was based on the complete set of genes in each strain, the
sample-specific behaviour was regained and the neural part improved by 0.035 PR-AUC
under the lineage holdout. Future efforts should be directed towards enhancing
detection and assembly of plasmids and other mobile genetic elements, which are
likely to harbour missing signals for the four challenging antibiotics. Future
holdout designs should also maintain a minimum number of resistant strains for each
antibiotic, particularly when the results are based upon only 38 cases of resistance
(as in the gentamicin score). Finally, the entire study should be performed with
another bacterial species to determine if the two-group phenomenon is unique to
\textit{Escherichia coli} or is a general pattern for genomic resistance prediction.
The main contribution of this work is not a single headline score. It is the finding
that genomic resistance prediction in \textit{E. coli} is not one problem but two,
and that the boundary between them is sharp, reproducible and specific to each
antibiotic. The model is what produced that finding, and the evaluation protocol is
what made it measurable.
